\documentclass[11pt]{article}

% Basic packages
\usepackage[T1]{fontenc}
\usepackage[utf8]{inputenc}
\usepackage{lmodern}
\usepackage[a4paper,margin=1in]{geometry}
\usepackage{amsmath,amssymb,amsthm,mathtools}
\usepackage{bm}
\usepackage{enumitem}
\usepackage{hyperref}
\usepackage{xcolor}
\usepackage{graphicx}
\usepackage{booktabs}
\usepackage{microtype}

% Theorem environments
\newtheorem{proposition}{Proposition}
\newtheorem{remark}{Remark}
\newtheorem{conjecture}{Conjecture}
\newtheorem{question}{Question}

% Shortcuts
\newcommand{\R}{\mathbb{R}}
\newcommand{\Pp}{\mathcal{P}_p}
\newcommand{\eps}{\varepsilon}
\newcommand{\iid}{\mathrel{\overset{\mathrm{i.i.d.}}{\sim}}}
\newcommand{\toP}{\xrightarrow{\mathbb P}}
\newcommand{\toD}{\xrightarrow{\mathcal D}}
\newcommand{\RightarrowW}{\Rightarrow}
\newcommand{\W}{W_p}
\newcommand{\permABC}{\textsc{permABC}}
\newcommand{\WABC}{\textsc{WABC}}

\title{Notes on New Theoretical Directions for \permABC}
\author{Working notes for discussion with co-authors}
\date{\today}

\begin{document}

\maketitle

\begin{abstract}
These notes collect a possible new theoretical angle for \permABC, centered on its relation to
Wasserstein ABC and on two asymptotic regimes that appear natural in the hierarchical setting:
large number of observed compartments $K \to \infty$, and large over-sampling size
$M \to \infty$. The goal is not to state definitive results yet, but to provide a clean framework
for discussion, isolate plausible statements, and identify the points that should be formalized in
future work.
\end{abstract}

\tableofcontents

\section{Motivation and purpose of these notes}

The purpose of this short note is to present a possible extension of the current \permABC
story in a form that can be discussed with co-authors before deciding whether any of it should
enter the paper.

The main ideas are the following.
\begin{enumerate}[label=(\arabic*)]
    \item \permABC can be viewed as a hierarchical analogue of \WABC at the level of
    empirical distributions of \emph{compartments}.
    \item This suggests a natural asymptotic regime when the number of compartments
    $K \to \infty$.
    \item The over-sampling bridge $M > K$ has a qualitatively different asymptotic behaviour:
    for large $M$, the global marginal tends to become prior-like, whereas the matched local
    variables become extreme-value / best-match objects, often nearly Dirac.
\end{enumerate}

The point of these notes is therefore to organize these observations into a coherent narrative and
make explicit which parts seem rigorous, which parts are currently heuristic, and which questions
could realistically be turned into formal results.

\input{sections/asymptotics.tex}

\section{Practical discussion points for co-authors}

At this stage, the most useful questions for discussion seem to be:
\begin{enumerate}[label=(\arabic*)]
    \item Should the $K \to \infty$ discussion appear in the main paper, the supplement, or only
    in internal notes for now?
    \item Do we want to state the large-$K$ argument only for the global parameter $\beta$,
    leaving the local parameters $\mu_{1:K}$ explicitly outside the main theorem?
    \item Do we want to include a formal proposition explaining why large-$M$ over-sampling is
    not an inferential target, but only an algorithmic bridge?
    \item Would it be worth adding a simple toy example in which the large-$M$ behaviour can be
    shown analytically, e.g. a deterministic or Gaussian local model?
    \item Is there a clean way to phrase the projected local variables as transport witnesses,
    selected latent explanations, or approximate Monge matches?
\end{enumerate}

\section{Suggested next steps}

A possible sequence of next steps would be:
\begin{enumerate}[label=(\arabic*)]
    \item Clean the notation so that the WABC comparison uses exactly the same symbols as the
    current paper.
    \item Decide which parts are intended as theorem/proposition material and which parts should
    remain discussion/conjecture.
    \item Build one toy model where the $M \to \infty$ phenomenon is explicit.
    \item If the co-authors agree, turn the large-$K$ part into a concise proposition linked to the
    WABC literature.
\end{enumerate}

\end{document}

